\documentclass[11pt,a4paper]{article}

% Paquetes basicos compatibles con pdfLaTeX/MiKTeX
\usepackage[utf8]{inputenc}
\usepackage[T1]{fontenc}
\usepackage[spanish]{babel}
\usepackage[a4paper,top=1.6cm,bottom=1.6cm,left=1.7cm,right=1.7cm]{geometry}
\usepackage{graphicx}
\usepackage{array}
\usepackage{tabularx}
\usepackage{hyperref}
\usepackage{xcolor}
\usepackage{float}

\setlength{\parindent}{0pt}
\setlength{\parskip}{4pt}
\renewcommand{\arraystretch}{1.18}
\hypersetup{colorlinks=true,urlcolor=blue,linkcolor=black}

\definecolor{azulOscuro}{RGB}{20,34,58}

\newcommand{\seccionazul}[1]{%
  \vspace{4pt}%
  \noindent\colorbox{azulOscuro}{%
    \parbox{\dimexpr\textwidth-2\fboxsep\relax}{\color{white}\bfseries #1}}%
  \vspace{4pt}%
}

\newcommand{\subpregunta}[1]{%
  \vspace{5pt}\noindent{\bfseries\color{azulOscuro} #1}\par
}

\newcommand{\imagen}[2]{%
\begin{figure}[H]
\centering
\IfFileExists{figuras/#1}{%
  \includegraphics[width=#2\textwidth,height=0.78\textheight,keepaspectratio]{figuras/#1}%
}{%
  \fbox{\parbox[c][4cm][c]{0.90\textwidth}{\centering
  \textbf{IMAGEN PENDIENTE}\\[4pt]
  Archivo: \texttt{figuras/#1}}}%
}
\end{figure}
}

\begin{document}

\begin{center}
{\Large\bfseries UNIVERSIDAD TÉCNICA ESTATAL DE QUEVEDO}\\[3pt]
{\bfseries Facultad de Ciencias de la Ingeniería - Carrera de Ingeniería de Software}\\[8pt]
{\large\bfseries PRUEBA PRÁCTICA - UNIDAD IV PARALELO B}\\
{\bfseries Validación, Gestión, Herramientas y Estándares}\\[8pt]
\end{center}

\begin{tabularx}{\textwidth}{>{\bfseries}l X >{\bfseries}l X}
Asignatura: & Ingeniería de Requisitos (ISR-401) & Docente: & Ing. Gleiston Guerrero, Mg.\\
Estudiante: & Anderson Jeampiere Naranjo Flores & Fecha: & 12 de Agosto de 2026\\
Modalidad: & Individual, en clase & Duración: & 120 minutos\\
Caso Evaluado: & Sistema de Reserva de Citas Médicas & Curso: & 4to Software B\\
\end{tabularx}

\vspace{6pt}
\textbf{URL del Repositorio Git Público (Criterio de Piso 1):}\\
\url{https://github.com/AnderJM3/Evaluacion-IR-Unidad-IV-Naranjo-Anderson}

\seccionazul{1. EVIDENCIA DE CUESTIONARIO SGA (COMPONENTE 1 - 30 PTS)}

\textbf{Resumen y Evaluación del Cuestionario rendido en el LMS (SGA):}\\
Nota obtenida: 18.70 / 20 pts.

\textbf{Captura 1: Resumen de Calificación SGA}
\imagen{captura_resumen_sga.png}{0.98}

\textbf{Captura 2: Evaluación / Intento Finalizado}
\imagen{captura_evaluacion_sga.png}{0.98}

\seccionazul{2. DESARROLLO DE ACTIVIDADES PRÁCTICAS (COMPONENTE 2 - 70 PTS)}

\subpregunta{P1. Modelo de Datos - Diagrama de Clases UML (8 pts)}
Modelo de dominio con atributos tipados, identificadores subrayados, operaciones clave y multiplicidades asociativas:

\imagen{p1_clases.png}{0.98}

\subpregunta{P2. Modelo Funcional - Diagrama de Actividades UML (8 pts)}
Flujo del proceso \textit{Agendar Cita} con decisión de cupo, caminos alternativos, convergencia y carriles de responsabilidad:

\imagen{p2_actividades.png}{0.98}

\subpregunta{P3. Modelo de Comportamiento - Máquina de Estados UML (8 pts)}
Ciclo de vida de una Cita con superestado para la transición cancelar y estado final de aceptación:

\imagen{p3_estados.png}{0.98}

\subpregunta{P4. Consistencia entre las Tres Perspectivas (7 pts)}
\textbf{Tabla de verificación cruzada y corrección de inconsistencia:}

{\footnotesize
\setlength{\tabcolsep}{4pt}
\begin{tabularx}{\textwidth}{|>{\raggedright\arraybackslash}p{0.24\textwidth}|>{\raggedright\arraybackslash}p{0.38\textwidth}|>{\raggedright\arraybackslash}X|}
\hline
\textbf{Evento (P3)} & \textbf{Función/Acción (P2)} & \textbf{Entidad / Atributo modificado (P1)}\\ \hline
\texttt{ev: verificarCupo} & Verificar Disponibilidad & Cupo.disponible\\ \hline
\texttt{ev: confirmar / agendar} & Confirmar Cita / Bloquear Cupo & Cita.estado, Cupo.disponible\\ \hline
\texttt{ev: notificarNoDisponibilidad} & Notificar No Disp. / Proponer Horario & Cita.estado\\ \hline
\texttt{ev: atender / cancelar / expirarHorario} & Acciones de ejecución/seguimiento & Cita.estado, Cupo.disponible\\ \hline
\end{tabularx}
}

\textbf{Defecto detectado y corrección:} El evento \texttt{ev:expirarHorario} (marcar no asistida) en P3 no tenía reflejo operacional directo en P1.

\textbf{Corrección:} Se añadió el método explícito \texttt{marcarNoAsistida()} en la clase Cita de P1 para mantener consistencia operacional.

\subpregunta{P5. Especificación de Requisitos con Esquema de Atributos (10 pts)}
Cuatro requisitos especificando métrica e ISO/IEC 25010:2023 para No Funcionales:

{\footnotesize
\setlength{\tabcolsep}{4pt}
\begin{tabularx}{\textwidth}{|>{\raggedright\arraybackslash}p{1.15cm}|>{\raggedright\arraybackslash}p{1.65cm}|>{\raggedright\arraybackslash}X|>{\raggedright\arraybackslash}p{2.8cm}|}
\hline
\textbf{ID} & \textbf{Tipo} & \textbf{Descripción / Métrica y Umbral (ISO 25010)} & \textbf{Prio / Est / Ver}\\ \hline
RF-01 & Funcional & El sistema debe permitir agendar una cita médica verificando la disponibilidad de cupos en la agenda del médico. & Alta / Alta / Prueba\\ \hline
RF-02 & Funcional & El sistema debe permitir cancelar una cita confirmada mientras esta no haya sido atendida aún. & Media / Alta / Prueba\\ \hline
RNF-01 & No Funcional & \textbf{ISO 25010: Eficiencia de Desempeño.} El tiempo de procesamiento para registrar y confirmar la reserva de cita debe ser $< 3$ segundos bajo 50 usuarios concurrentes. & Alta / Media / Inspección\\ \hline
RNF-02 & No Funcional & \textbf{ISO 25010: Seguridad (Protección de datos).} Los datos personales y de contacto del paciente (cédula, correo) deben cifrarse mediante AES-256 en reposo y tránsito. & Alta / Alta / Análisis\\ \hline
\end{tabularx}
}

\subpregunta{P6. Priorización MoSCoW (5 pts)}

{\footnotesize
\setlength{\tabcolsep}{4pt}
\begin{tabularx}{\textwidth}{|>{\raggedright\arraybackslash}p{1.35cm}|>{\raggedright\arraybackslash}p{1.9cm}|>{\raggedright\arraybackslash}X|}
\hline
\textbf{Req.} & \textbf{Categoría} & \textbf{Justificación breve (Valor vs. Viabilidad)}\\ \hline
RF-01 & Must & Es la funcionalidad núcleo sin la cual el sistema no tiene razón operativa de ser.\\ \hline
RNF-02 & Must & Garantiza el cumplimiento normativo e legal obligatorio en protección de datos de salud.\\ \hline
RNF-01 & Should & Aporta agilidad a la atención del paciente, técnicamente factible con la arquitectura planificada.\\ \hline
RF-02 & Could & Aporta flexibilidad al paciente, pero las cancelaciones pueden gestionarse vía recepción en fase 1.\\ \hline
\end{tabularx}
}

\subpregunta{P7. Validación por Inspección - Lista de Comprobación ISO/IEC/IEEE 29148 (8 pts)}

\textbf{1. Registro de Defectos Detectados (Sin corregir):}
\begin{itemize}
\item \textbf{Defecto RF-01:} Ambigüedad al no precisar qué sucede con el cupo si la cita es rechazada por falta de horario.
\item \textbf{Defecto RNF-01:} Incompleto al no indicar las condiciones del entorno/servidor donde se mide el umbral de 3s.
\end{itemize}

\textbf{2. Retrabajo (Requisitos Corregidos):}
\begin{itemize}
\item \textbf{RF-01 (Corregido):} El sistema debe registrar una cita validando el cupo; si está disponible, bloqueará el cupo; si no, notificará la no disponibilidad y sugerirá alternativas.
\item \textbf{RNF-01 (Corregido):} El sistema debe procesar el registro de reserva en $< 3$s bajo 50 usuarios concurrentes en el entorno de pruebas homologado.
\end{itemize}

\subpregunta{P8. Pruebas de Aceptación Trazadas (6 pts)}

{\scriptsize
\setlength{\tabcolsep}{3pt}
\begin{tabularx}{0.97\textwidth}{|>{\raggedright\arraybackslash}p{1.15cm}|>{\raggedright\arraybackslash}p{1.2cm}|>{\raggedright\arraybackslash}p{1.55cm}|>{\raggedright\arraybackslash}X|>{\raggedright\arraybackslash}X|}
\hline
\textbf{ID Prueba} & \textbf{Req.} & \textbf{Tipo} & \textbf{Entradas de Prueba} & \textbf{Criterio de Éxito}\\ \hline
TC-01 & RF-01 & Funcional & Paciente con cédula válida, Médico A con cupo disponible a las 10:00. & Cita agendada en estado Confirmada y cupo marcado como no disponible.\\ \hline
TC-02 & RF-02 & Funcional & ID de cita en estado Confirmada, solicitud de cancelación recibida. & Cita cambia a estado Cancelada y el cupo se libera en la agenda.\\ \hline
TC-03 & RNF-01 & No Funcional & 50 peticiones simultáneas de reserva de citas mediante JMeter. & El 95\% de las transacciones responde en menos de 3 segundos.\\ \hline
TC-04 & RNF-02 & No Funcional & Inspección de la BD de citas y captura de tráfico de red. & Cédula y correo almacenados con AES-256 y canal protegido con TLS 1.3.\\ \hline
\end{tabularx}
}

\subpregunta{P9. Matriz de Trazabilidad (6 pts)}

{\scriptsize
\setlength{\tabcolsep}{2.5pt}
\begin{tabularx}{0.97\textwidth}{|>{\centering\arraybackslash}p{1.15cm}|>{\raggedright\arraybackslash}p{2.15cm}|>{\raggedright\arraybackslash}X|>{\centering\arraybackslash}p{1.15cm}|>{\centering\arraybackslash}p{1.55cm}|>{\centering\arraybackslash}p{1.40cm}|}
\hline
\textbf{Req.} & \textbf{Fuente (Atrás)} & \textbf{Modelos UML (Adelante)} & \textbf{Prueba} & \textbf{Dep. horizontal} & \textbf{Huérfano}\\ \hline
RF-01 & Paciente / Clínica & P1 (Cita, Cupo), P2, P3 & TC-01 & RNF-01 & No\\ \hline
RF-02 & Clínica / Políticas & P1 (Cita), P3 (Superestado) & TC-02 & RF-01 & No\\ \hline
RNF-01 & Calidad / SLA & P2 (Proceso Agendar) & TC-03 & Ninguna & No\\ \hline
RNF-02 & Seguridad / Ley & P1 (Paciente) & TC-04 & Ninguna & No\\ \hline
\end{tabularx}
}

\seccionazul{3. MARCO NORMATIVO DE REFERENCIA}

\begin{enumerate}
\item Fagan, M. E. (1976). \textit{Design and code inspections to reduce errors in program development}. IBM Systems Journal.
\item ISO/IEC 25010:2023. \textit{Systems and software engineering - Systems and software quality requirements and evaluation (SQuaRE)}.
\item ISO/IEC/IEEE 29148:2018. \textit{Systems and software engineering - Life cycle processes - Requirements engineering}.
\item Object Management Group (OMG). (2017). \textit{Unified Modeling Language (UML), version 2.5.1}.
\item Pohl, K. (2010). \textit{Requirements Engineering: Fundamentals, Principles, and Techniques}. Springer.
\item IEEE Computer Society. (2024). \textit{SWEBOK Guide v4.0}.
\end{enumerate}

\end{document}
